\scanpage{115}
\mspage{705}

\[
\tag{267}
\left\{
\begin{gathered}
\unsure{\rho'}=\zeta\rho,\ \sigma'=\eta\sigma
\quad(\zeta,\eta\in\Gamma\underline M)\quad\text{such that we have}
\\
\rho'{}^3=\sigma'{}^2=\omega,\qquad
\sigma'(\rho')=\rho'{}^{-1}.
\end{gathered}\right.
\]

(We \unsure{note|shall note} that a splitting must transform $\omega$,
which operates trivially on $\underline M$ and is central
in $\underline H^+$, into a central \deleted{\unsure{element}}
\inserted{section}, \unsure{now}
\deleted{\unsure{the inverse}} \unsure{see|here is} \unsure{why} we have
\[
\tag{268}
\left\{\begin{aligned}
\underline M\cap\underline{\mathrm{Cent}}\bigl(\underline{\mathfrak S}^+\bigr)
&=\underline M^{\mathfrak S_3}
=\deleted{\unsure{\{0\}}}\,\unsure{\{1\}|\{0\}},\quad\text{hence}
\\
\underline{\mathrm{Cent}}\bigl(\underline{\mathfrak S}^+\bigr)
&\simeq\underline\mu_{\underline{\mathfrak S}}
=\{1,\omega\}_{\underline{\mathfrak S}}
\end{aligned}\right.
\]
--- which shows that a splitting must transform $\omega$ into $\omega$,
whence \unsure{easily} the characterization (267) of the $\rho',\sigma'$
\unsure{corresponding|corresponding} to a splitting).
We shall have from this
\[
\left\{\begin{aligned}
\rho'{}^3&=(\zeta\rho)^3
=\zeta\rho(\zeta)\rho^2(\zeta)\underbrace{\rho^3}_{\omega}
=\zeta\rho(\zeta)\rho^2(\zeta)\omega
\\
\sigma'{}^2&=(\eta\sigma)^2
=\eta\sigma(\eta)\underbrace{\sigma^2}_{\omega}
=\eta\sigma(\eta)\omega
\end{aligned}\right.
\]
thus the \deleted{two} equations $\rho'{}^3=\omega$ and
$\sigma'{}^2=\omega$ are equivalent respectively to
\[
\left\{\begin{aligned}
\zeta\rho(\zeta)\rho^2(\zeta)&=1
&\text{i.e.}\quad(1+\rho_M+\rho_M^2)(\zeta)&=0
\\
\eta\sigma(\eta)&=1
&\text{i.e.}\quad(1+\sigma_M)(\eta)&=0
\end{aligned}\right.
\]
i.e. this reduces to equations (245) (246),
and we have seen that the \deleted{general solution is}
\inserted{\deleted{first} condition}
on $\zeta$ is automatically satisfied, that
on $\eta$ amounts to:
\[
(247)\qquad
\eta=b\delta=b(\lambda_1-\lambda_0),\quad
\unsure{\lambda|b}\in\Gamma\mathbb G_a.
\]
The condition $\sigma'(\rho')=\rho'{}^{-1}$ is made explicit on the other
hand as on page~$700'$, we come back

\margin{114}

\begin{center}
\small
114\\
\texttt{https://grothendieck.umontpellier.fr}\\
Université de Montpellier, UMR CNRS 5149, Case courrier 051, Place Eugène Bataillon\\
34095 Montpellier cedex 5 - France
\end{center}

\scanpage{116}

to equation (249), which is made explicit by (250).
This shows that there exists a splitting
even over $S_0=\mathrm{Spec}\,\mathbb Z$
--- which amounts to a splitting \unsure{of the} discrete group
$\mathfrak S^+$ extension of
$H^+=\mathfrak S^+/M$
\inserted{$\simeq\mathfrak S_0/\Pi_0$}
(by $M\simeq(\Pi_0)_{\mathrm{ab}}$ \unsure{)}
--- and that the splittings in question form a
torsor under $\underline K$, where $K$\inserted{$\simeq\mathbb Z^2$}
is defined in (253), and in fact \deleted{\unsure{$K$}}
we have an isomorphism $M\simeq K$ (254), i.e.
$\underline M\simeq\underline K$, \unsure{corresponding}
to the operation of $\underline M$ on the \unsure{scheme}
of splittings, by conjugation. Thus we find

\uline{Proposition} (Over an arbitrary base $S$).

The extension $\underline{\mathfrak S}^+$ of
$\underline H^+$ by $\underline M$ is split ---
two splittings are conjugate by a unique section
of $\underline M$ --- thus the map \deleted{of}
``passage to the quotient by $\mu_{\mathfrak S}$'' defines a
bijection from the \unsure{set} of splittings of
$\underline{\mathfrak S}^+$ over $\underline H^+$,
onto the \unsure{set} of splittings of
$\underline{\mathfrak S}^+/\mu_{\underline{\mathfrak S}}$
over $\underline H^+/\underline\mu_{\underline H^+}
\simeq\underline{\mathfrak S}_3$.

\margin{115}

\begin{center}
\small
115\\
\texttt{https://grothendieck.umontpellier.fr}\\
Université de Montpellier, UMR CNRS 5149, Case courrier 051, Place Eugène Bataillon\\
34095 Montpellier cedex 5 - France
\end{center}

\scanpage{117}
\mspage{706}

Let us now study the centralizer in
$\underline G_{\rho,\sigma}$ of a section subgroup
over $\underline H^+$, \unsure{say} denoted
$\underline H(\rho',\sigma')$. We must now have
\[
\tag{269}
\underline{\mathrm{Cent}}_{\underline G_{\rho,\sigma}}
\bigl(\underline H(\rho',\sigma')\bigr)
\subset
\begin{gathered}
\mathbb O_{12}^*\cdot\bigl(\underline M\times
\underbrace{\mu_{\underline{\mathfrak S}}}_{=\{1,\omega\}_{S_0}}\bigr)
\\[-2pt]
\Vert
\\[-2pt]
\left\{\lambda\omega^i u_{\underline\mu}\ \middle|\
\begin{gathered}
\lambda\in\Gamma\underline M\\
\underline\mu\in\Gamma\mathbb O_{12}^*\\
i\in\Gamma(\unsure{\mathbb Z/2\mathbb Z})_{S_0}
\end{gathered}\right\}
\end{gathered}
\]
Moreover $\lambda\omega^i u_{\underline\mu}$ centralizes
$\underline H(\rho',\sigma')$ iff $\lambda u_{\underline\mu}$
centralizes. Commutation with $\rho'$ is again made explicit
by condition (260), the calculation on page~703 applying
\unsure{ne varietur}. For the condition of commutation with
$\sigma'$, the calculation on page~$703'$ must be modified
for we can no longer write $\omega^{\nu_2}=1$,
the condition is written
\[
\omega^{\nu_2}\cdot\lambda\eta^{\unsure{\mu}}
\sigma(\lambda^{-1})\sigma=\eta\sigma
\quad\text{i.e.}\quad
\omega^{\nu_2}\lambda\eta^\mu\sigma(\lambda^{-1})=\eta
\]
and is equivalent to the conjunction of the condition
\[
\tag{270}
\nu_2(\underline\mu)=0\quad\text{i.e.}\quad
\underline\mu\in\Gamma\mathbb O_{12}^{*\prime\prime},
\]
and of relation (261). In the presence of
$\nu_2(\underline\mu)=0$, the calculations on pages~$703'$, 704
apply, \unsure{finally} we find:

\uline{Corollary} For every splitting of
$\underline{\mathfrak S}^+$ over
$\underline H^+$, \deleted{the commutant}
\inserted{\deleted{\illegible{1}} let $\underline T$ be the centralizer
\deleted{\unsure{$T$}}} in $\underline G_{\rho,\sigma}$
of the section subgroup $\underline H(\rho',\sigma')$,
\deleted{we}

\margin{116}

\begin{center}
\small
116\\
\texttt{https://grothendieck.umontpellier.fr}\\
Université de Montpellier, UMR CNRS 5149, Case courrier 051, Place Eugène Bataillon\\
34095 Montpellier cedex 5 - France
\end{center}

\scanpage{118}

\deleted{\unsure{a section}}\inserted{We then have} an exact sequence
\[
\tag{271}
\diagram{\begin{tikzcd}[column sep=large]
1 \arrow[r] & \underline\mu_{\unsure{\mathfrak S|G}}
\arrow[r,hook,"{\deleted{\unsure{G}}}" description]
& \underline T \arrow[r] & \mathbb O_{12}^{*\prime\prime}
\arrow[r,no head] & 1
\end{tikzcd}}
\]
and for $\underline T^{\unsure{\natural|\sharp}}
=\underline T\cap\underline G_{\rho,\sigma}^{\unsure{\natural|\sharp}}$,
we have
\[
\tag{272}
\underline T^{\unsure{\natural|\sharp}}
\simeq\mathbb O_{12}^{*\prime\prime},
\]
i.e. $\underline T^{\unsure{\natural|\sharp}}$ is a partial
section of $\underline G_{\rho,\sigma}^{\unsure{\natural|\sharp}}$
over $\mathbb O_{12}^{*\prime\prime}$.

\uline{Corollary} The extension $\underline G_{\rho,\sigma}$ of
$\mathbb O_{12}^*\times\underline H^+$ by $\underline M$
is not split over any nonempty base $S$.
The extension $\underline G_{\rho,\sigma}''$ of
$\mathbb O_{12}^{*\prime\prime}\times\underline H^+$ by
$\underline M$ is split over every base $S$,
and two splittings are conjugate by a unique section of
$\underline M$. The natural map gives a bijection
of the \unsure{set} of splittings of $\underline G_{\rho,\sigma}''$
over $\mathbb O_{12}^{*\prime\prime}\times\underline H^+$, with
the \unsure{set} of splittings of
$\underline{\mathfrak S}^+$ over $\underline H^+$.

\uline{Corollary} The extension $G_{\rho,\sigma}$ of
$\widehat{\mathbb Z}^{*}\times H^+$ by $M^\wedge$
is not split. That of $G_{\rho,\sigma}''$ of
$\widehat{\mathbb Z}^{*\prime\prime}\times H^+$ by
$M^\wedge$ is split, and two splittings are
conjugate by a unique element \unsure{of} $\widehat M$.
The natural map defines a bijection of the \unsure{set}
of splittings of $G_{\rho,\sigma}''$ over
$\widehat{\mathbb Z}^{*}\times H^+$ with the \unsure{set}
of splittings of \deleted{\unsure{$G_\rho$}}
$\mathfrak S^\wedge$ over $H^+$.

\margin{117}

\begin{center}
\small
117\\
\texttt{https://grothendieck.umontpellier.fr}\\
Université de Montpellier, UMR CNRS 5149, Case courrier 051, Place Eugène Bataillon\\
34095 Montpellier cedex 5 - France
\end{center}

\scanpage{119}
\mspage{707}

Let us make explicit a splitting of the extension
$\mathfrak S^+$ of $H^+$ by $M$
--- corresponding to a splitting of
$\underline{\mathfrak S}^+$ of $\underline H^+$
by $\underline M$ --- , via a particular solution of the equation
\[
\tag{273}
a_0-2a_1+3b-1=0
\]
taking
\[
a_1=b=0\qquad\text{hence}\qquad a_0=1
\]
i.e.
\[
\tag{274}
\left\{\begin{aligned}
\zeta&=\lambda_0,\ \eta=0\quad\text{(additive notation)}\\
\rho'&=\lambda_0\rho,\quad\unsure{\sigma'}=\unsure{\sigma}
\end{aligned}\right.
\]
\deleted{the corresponding splitting}
\inserted{the \unsure{three} formulas give a splitting of
$\mathfrak S^+/\mu_{\mathfrak S}$}
over $\mathfrak S_3
\simeq H^+/\mu_{H^+}$.
Let us make explicit the \unsure{corresponding}
\deleted{splittings}\inserted{liftings} of $\mathbb O_{12}^*$ into
$\underline G_{\rho,\sigma}/\mu_G$ resp. of
$\mathbb O_{12}^{*\prime\prime}$ into $\underline G_{\rho,\sigma}$\unsure{;},
\[
\tag{275}
\left\{\begin{aligned}
u_{\rho',\sigma'}:\mathbb O_{12}^*&\longrightarrow
\underline G_{\rho,\sigma}/\mu_G
\quad\text{resp.}\quad\mathbb O_{12}^{*\prime\prime}
\mathrel{\text{---}}\underline G_{\rho,\sigma}
\\
\underline\mu&\longmapsto\lambda_\infty^{\nu}u_{\underline\mu}
\qquad(\nu=\nu(\underline\mu))
\end{aligned}\right.
\]
where $u_{\underline\mu}$ defined in (151).
(N.B.\unsure{\S} (275) is contained in proposition page~704
(265) and (264), which give
\[
u_{\rho',\sigma'}=\lambda u_{\underline\mu}\qquad
\lambda=x\lambda_0+y\lambda_1\qquad x=y=-\nu
\]
\[
\text{hence}\qquad\lambda=-\nu(\lambda_0+\lambda_1)=\nu\lambda_\infty.
\]

We shall thus have, putting
\[
u'_{\underline\mu}=u_{\rho',\sigma'}(\underline\mu)
\]
\[
\tag{276}
\left\{\begin{aligned}
u'_{\underline\mu}(\rho')&=\rho',\quad
u'_{\underline\mu}(\sigma')=\sigma'\\
u'_{\underline\mu}(\lambda)&=\mu\lambda\quad
(\lambda\in\Gamma\underline M)
\quad\inserted{\text{where}}\ (\mu=\unsure{\chi}(\underline\mu))
\end{aligned}\right\}
\quad\begin{gathered}
\text{(which suffices}\\
\text{to \unsure{condition|characterize}}\ u'_{\underline\mu})
\end{gathered}
\]

\margin{118}

\begin{center}
\small
118\\
\texttt{https://grothendieck.umontpellier.fr}\\
Université de Montpellier, UMR CNRS 5149, Case courrier 051, Place Eugène Bataillon\\
34095 Montpellier cedex 5 - France
\end{center}

\scanpage{120}

and we also deduce from this
\[
\tag{277}
\left\{
\begin{aligned}
u'_{\underline{\mu}}(\rho)&=\lambda_0^{-2\nu}\rho\\
u'_{\unsure{\mu|\underline{\mu}}}(\sigma)&=\omega^{\nu_2}\sigma
\qquad (=\sigma\text{ in }\underline G_{\rho,\sigma}/\underline{\mu}_G\ \ldots).
\end{aligned}
\right.
\]

\margin{119}

\begin{center}
\small
119\\
\texttt{https://grothendieck.umontpellier.fr}\\
Université de Montpellier, UMR CNRS 5149, Case courrier 051, Place Eugène Bataillon\\
34095 Montpellier cedex 5 - France
\end{center}

\scanpage{121}
\mspage{708}

13\textsuperscript{o}) \uline{Return to the Fermat tower}.

Let us begin by bringing out a general topossic notion, which \unsure{was} a little
muddled between the lines in 11\textsuperscript{o}).
Let $\Gamma$ (here $\Gamma_{\mathbb Q}\times
\mathfrak S_3\deleted{\unsure{)|\}}}=N_{0,3}$) be operating on
a topos $\mathcal X$ (here
$\deleted{\unsure{U}}\inserted{\unsure{\mathcal U}_{0,3\overline{\mathbb Q}}}
\simeq M'_{0,4\overline{\mathbb Q}}$) connected
\inserted{locally connected}, \deleted{\unsure{and connected}} equipped with a ``point'' (geometric)
$x_0$ (here: the point $-\overline j$ with values in $\overline{\mathbb Q}$), let
\margin{Discrete context ($\mathcal X$ locally simply connected) or profinite,
\unsure{relatively}\illegible{1} \unsure{reasonable}\ldots}
\[
\tag{278}
\left\{
\begin{aligned}
\Pi&=\Pi_1(\mathcal X,x_0),\qquad
\mathcal M=\Pi_1(\mathcal X,\Gamma;x_0)\\
M&=\Pi_{\mathrm{ab}}=\Pi/[\Pi,\Pi],\qquad
\deleted{\unsure{\mathfrak S}}\inserted{\widetilde{\mathcal M}}
=\mathcal M/[\Pi,\Pi],\qquad \deleted{\text{hence}}
\end{aligned}
\right.
\]
hence exact sequences and homomorphisms of exact sequences
\[
\tag{279}
\diagram{\begin{tikzcd}[column sep=large,row sep=large]
1 \arrow[r,dash] & \Pi \arrow[r] \arrow[d,"\text{epimorphism}"]
& \mathcal M \arrow[r] \arrow[d,"\text{epimorphism}"]
& \Gamma \arrow[r,dash] \arrow[d,"\unsure{\sim|\wr}"] & 1\\
1 \arrow[r,dash] & M \arrow[r,dash]
& \deleted{\unsure{\mathfrak S}}\inserted{\widetilde{\mathcal M}} \arrow[r]
& \Gamma \arrow[r,dash] & 1
\end{tikzcd}}\ .
\]

\deleted{Let us consider}
\inserted{(Instead of taking
$M=\Pi/\unsure{(\Pi,\Pi)|[\Pi,\Pi]}\inserted{=\Pi_{\mathrm{ab}}}$,}
we \unsure{would have|had} been able to take a quotient of $\Pi_{\mathrm{ab}}$ invariant under
the action of $\Gamma$, but no matter\ldots).
Let us consider ``the universal abelian covering'' \inserted{$\mathcal X'_0$} of
$\mathcal X$ at $x_0$, i.e. the quotient by $[\Pi,\Pi]$ of the universal covering of
$\mathcal X$ at $x_0$ --- it is thus a principal covering of $\mathcal X$, with
group $M$. We

\margin{120}

\begin{center}
\small
120\\
\texttt{https://grothendieck.umontpellier.fr}\\
Université de Montpellier, UMR CNRS 5149, Case courrier 051, Place Eugène Bataillon\\
34095 Montpellier cedex 5 - France
\end{center}